

Nous allons appliquer ces définitions aux @@M0@@-anneaux augmentés @@M1@@ et @@M2@@, les anneaux d'augmentation @@M3@@ et @@M4@@, étant binômiaux (V 2.8). Pour simplifier les notations, nous poserons, pour tout schéma @@M5@@, et tout @@M6@@:
@@M7@@
Nous considérerons pour toute algèbre graduée @@M8@@ le groupe @@M9@@ comme contenu dans @@M10@@. Enfin, si on a un homomorphisme @@M11@@, additif, et compatible aux graduations à un décalage près, on notera encore @@M12@@ l'homomorphisme @@M13@@ qu'on en déduit.

\textbf{3.4.} Soit @@M14@@ un schéma. Pour tout @@M15@@-Module localement libre de type fini @@M16@@, on note @@M17@@ la classe dans @@M18@@ du dual @@M19@@ de @@M20@@. Il est clair que la fonction qui à @@M21@@ associe @@M22@@ est une fonction additive sur la catégorie des @@M23@@-Modules localement libres de type fini, et par suite définit un endomorphisme de @@M24@@, qu'on notera @@M25@@.

\textbf{Proposition 3.5.} \uline{Soit} @@M26@@ \uline{un schéma quasi-compact}. \uline{Alors pour tout} @@M27@@ \uline{et tout} @@M28@@, \uline{on a}:
@@M29@@

La relation à démontrer peut encore s'écrire :
@@M30@@
en posant : @@M31@@.

Par additivité, on peut donc supposer que @@M32@@ est la classe d'un @@M33@@-Module localement libre de type fini @@M34@@. De plus, quitte à remplacer @@M35@@ par le schéma des drapeaux de @@M36@@, ce qui donne une injection sur les
